\section*{Chapitre \ref{ch:m2:the-fx-options-market} --- Le marché des options de change}

\begin{solution}{exo:m2:the-fx-options-market:1}
$7 + 0.2 - 0.25 = 6.95\%$ pour le call sur l'euro à 25 delta et $7 + 0.2 + 0.25 = 7.45\%$ pour
le put~: le put sur l'euro est plus cher, d'un demi-point, soit la valeur du \omterm{def:m2:the-fx-options-market:rr}{risk reversal}.
\end{solution}

\begin{solution}{exo:m2:the-fx-options-market:2}
\omterm{def:m2:the-fx-options-market:delta}{Delta spot} usuel pour EURUSD (prime en dollars, devise de cotation)~; \omterm{def:m2:the-fx-options-market:delta}{delta ajusté de la prime} pour USDJPY
(prime en dollars, devise de base)~; USDBRL en \omterm{def:m2:the-fx-options-market:delta}{delta forward} et ajusté
de la prime (une paire de marchés émergents, prime en dollars). La prime payée
dans la devise de base est une position dans cette devise, ce qui modifie la couverture et donc le
delta qui définit chaque strike coté.
\end{solution}

\begin{solution}{exo:m2:the-fx-options-market:3}
$600\,000 - 73\,669 = 526\,331$ euros~: un \omterm{def:m2:the-fx-options-market:delta}{delta ajusté de la prime} de 52.63~\%.
\end{solution}

\begin{solution}{exo:m2:the-fx-options-market:4}
\omterm{def:m2:the-fx-options-market:delta}{Deltas spot} usuels~: call 1.17904, put 1.12357~; \omterm{def:m2:the-fx-options-market:delta}{deltas forward}~: 1.17920 et
1.12341.
\end{solution}

\begin{solution}{exo:m2:the-fx-options-market:5}
Avec des deltas usuels, le delta du call $N(d_1)$ égale celui du put $N(-d_1)$ lorsque
$d_1 = 0$, soit $K = Fe^{\sigma^2T/2}$, au-dessus du forward. Le \omterm{def:m2:the-fx-options-market:delta}{delta ajusté de la prime}
utilise $N(d_2)$ multiplié par $K/F$, et la neutralité exige $d_2 = 0$, soit
$K = Fe^{-\sigma^2T/2}$, en dessous du forward.
\end{solution}

\begin{solution}{exo:m2:the-fx-options-market:6}
Il ne verse rien dans les trajectoires où le spot touche la barrière avant l'échéance, certaines
d'entre elles auraient fini dans la monnaie. L'écart est grand lorsque la barrière
est proche du spot, la volatilité élevée et l'échéance lointaine, puisque le touché est
alors probable.
\end{solution}

\begin{solution}{exo:m2:the-fx-options-market:7}
Avec des \omterm{def:m2:the-fx-options-market:delta}{deltas forward}, le strike du call est 160.90 et celui du put 150.65, contre 160.85 et
150.71 avec les \omterm{def:m2:the-fx-options-market:delta}{deltas spot} usuels, et 160.68 et 150.52 avec les \omterm{def:m2:the-fx-options-market:delta}{deltas spot}
ajustés de la prime.
\end{solution}

\begin{solution}{exo:m2:the-fx-options-market:8}
Un \omterm{def:m2:the-fx-options-market:rr}{risk reversal} est un prix, non une prévision~: les puts en dollars sont plus chers que les calls
car davantage d'acheteurs veulent se protéger contre une baisse, et les dealers qui la vendent
facturent le risque d'un mouvement brutal dans ce sens. Il signifie que le marché paie plus pour
se protéger contre une baisse du dollar~; c'est le forward, et non le smile, qui porte
l'espérance du taux du marché sous la mesure de valorisation.
\end{solution}

\begin{solution}{pb:m2:the-fx-options-market:1}
\textbf{1.} 183.8 pips de notionnel (0.01838 dollar par euro).
\textbf{2.} 165.2 pips~; l'entreprise économise 18.6 pips, soit environ 930\,000~USD sur
500~millions d'euros.
\textbf{3.} 0.657~; le dealer détient environ 328.7~millions d'euros.
\textbf{4.} Près de la barrière, la valeur du knock-out tombe à zéro bien plus vite que
celle de l'option vanille lorsque le spot baisse, donc sa valeur est plus sensible au spot.
\textbf{5.} 0.631 juste au-dessus de 1.1200, contre 0.298 pour l'option vanille.
\textbf{6.} Environ 315.5~millions d'euros, en position longue en euros.
\textbf{7.} Vendre la totalité, soit 315.5~millions d'euros, au touché.
\textbf{8.} Cela ajoute une vente importante au moment où le spot traverse la barrière à la baisse,
ce qui fait encore baisser le spot.
\textbf{9.} Acheter des euros près de 1.1200, pour éviter que la barrière ne soit touchée et que l'option
reste active.
\textbf{10.} S'ils savent qu'une couverture importante se dénoue là, pousser le spot jusqu'à ce niveau peut
déclencher la vente et un mouvement dont ils peuvent tirer profit.
\textbf{11.} En vendant par avance une partie de la couverture à mesure que le spot approche, ou en
achetant des options vanille qui compensent le saut de delta de la barrière.
\textbf{12.} Elle sacrifie la précision de la couverture~: si le spot rebrousse chemin, le dealer
reste vendeur du delta vendu tôt~; les options coûtent leur prime.
\textbf{13.} Le gamma près d'une barrière knock-out est élevé et change de signe de part et d'autre~:
de petits mouvements exigent des transactions de rééquilibrage importantes, chacune payant la fourchette.
\textbf{14.} Avec un smile, la volatilité à la barrière (celle du put sur l'euro) est plus élevée,
la probabilité de toucher la barrière plus forte, et le knock-out moins cher qu'avec la
volatilité plate.
\textbf{15.} Que la protection disparaît précisément lorsque le spot traverse 1.1200,
et que sa décote est le prix de ce risque.
\textbf{16.} Couvrir une exposition est légitime~; opérer pour empêcher ou provoquer un
touché afin d'en tirer profit, en utilisant des informations clients, ne l'est pas, et relève
des règles contre les manipulations et du \omterm{def:m2:fx-spot-microstructure:code}{FX Global Code}.
\textbf{17.} Les clients choisissent des niveaux ronds, faciles à retenir et à surveiller.
\textbf{18.} Les niveaux de barrière et les montants des clients, qui révèlent où d'autres pourraient pousser
le marché.
\textbf{19.} Résultat à retenir~: \emph{la défense de barrière} fixe la vente du dealer au
touché, soit environ 315.5~millions d'euros, un delta de 0.631 sur 500~millions d'euros, le tout à
1.1200.
\textbf{20.} Parce que les couvertures des options qui disparaissent ou apparaissent à un niveau doivent être
dénouées ou mises en place à ce niveau, en volume et immédiatement.
\end{solution}

\begin{solution}{iq:m2:the-fx-options-market:1}
La volatilité à la monnaie est le niveau~; le \omterm{def:m2:the-fx-options-market:rr}{risk reversal} à 25 delta, soit la volatilité du call
moins celle du put, est le skew~: il indique de quel côté le marché paie le plus pour s'assurer~; le
\omterm{def:m2:the-fx-options-market:rr}{butterfly}, soit la moyenne des ailes par rapport à la monnaie, est la courbure~: elle indique comment l'épaisseur des queues
est valorisée. Ensemble, ils donnent trois points du smile.
\iqlookfor{niveau, skew et courbure.}
\end{solution}

\begin{solution}{iq:m2:the-fx-options-market:2}
Le delta corrigé d'une prime payée dans la devise de base, qui est elle-même une
position dans cette devise~: $\phi(K/F)N(\phi d_2)$ en termes de forward, ou actualisé
pour le spot. C'est la convention pour les paires dont la prime est payée dans la devise de base
comme USDJPY ou EURJPY.
\iqlookfor{pourquoi la prime modifie la couverture, et pour quelles paires.}
\end{solution}

\begin{solution}{iq:m2:the-fx-options-market:3}
Convertir les cotations en volatilité du put, par la formule simplifiée ou par une
calibration sur le \omterm{def:m2:the-fx-options-market:rr}{strangle} de marché~; puis résoudre le strike auquel un put avec
cette volatilité a un delta de $-0.25$ selon la convention de la paire~: forme fermée pour
les deltas usuels, une recherche monotone de racine pour les \omterm{def:m2:the-fx-options-market:delta}{deltas ajustés de la prime}.
\iqlookfor{l'ordre~: d'abord la volatilité, puis le strike, dans le bon delta.}
\end{solution}

\begin{solution}{iq:m2:the-fx-options-market:4}
Estimer le delta qui sera libéré au touché et le gamma en chemin~;
vendre une partie de la couverture à l'avance ou acheter des options qui compensent le saut, dans des limites~;
garder la barrière confidentielle~; ne pas opérer pour provoquer ou empêcher le touché~; si
le touché survient, exécuter le dénouement comme prévu plutôt que de courir après le marché.
\iqlookfor{préparation, et non manipulation.}
\end{solution}

\begin{solution}{iq:m2:the-fx-options-market:5}
Parce que le straddle à ce strike n'a pas de delta, il se négocie comme un instrument
de pure volatilité et ne nécessite aucune couverture sur la transaction~; le strike forward laisserait
un petit delta. La convention reflète ce à quoi servent les options à la monnaie,
à savoir négocier la volatilité.
\iqlookfor{la neutralité en delta comme objectif.}
\end{solution}

\begin{solution}{iq:m2:the-fx-options-market:6}
Stocker pour chaque paire ses conventions (type de delta par maturité, définition de l'ATM, devise de la
prime)~; intégrer les ATM, RR et BF par maturité avec horodatage~; calibrer un smile par
échéance, exactement sur les cotations~; interpoler dans le temps la variance totale~; servir
les volatilités par strike ou par delta, avec la convention énoncée~; versionner chaque surface~;
vérifier l'absence d'arbitrage (calendaire et \omterm{def:m2:the-fx-options-market:rr}{butterfly}) ainsi que les cotations périmées~; tester contre
des cas construits à la main.
\iqlookfor{conventions comme données, calibration exacte, contrôles d'arbitrage.}
\end{solution}
